\scanpage{287}
\mspage{15}

\[
 \operatorname{Hom}_{\mathrm{Ens}}(\pi_0(X),I)\simeq\operatorname{Hom}_W(X,I_W).
 \tag*{(2.8.2)}
\]

\deleted{Quand $W$ satisfait l'axiome $\operatorname{Loc}(4a)$ des sommes \unsure{d'ordre|directes}, alors} \inserted{Le foncteur} $\pi_0$ est également \uline{adjoint à droite} du \uline{même} foncteur $I\mapsto I_W$, i.e. on a une bijection fonctorielle

\[
 \operatorname{Hom}_{\mathrm{Ens}}(I,\pi_0(X))\simeq\operatorname{Hom}_W(I_W,X)\unsure{;|.}
 \tag*{(2.8.3)}
\]

Prenant $I=e$ (ensemble ponctuel), on trouve le \deleted{corollaire :}

\uline{Corollaire} 2.9. \deleted{Moyennant $\operatorname{Loc}(4a)$,} On a \deleted{\unsure{la bijec|la bijection}} l'isom fonctoriel en $X$ de $\mathrm{Hot}_W$

\[
 \boxed{\operatorname{Hom}_W(e_W,X)\simeq\pi_0(X)}. \tag*{(2.8.4)}
\]

\uline{Commentaire} 2.10. La formule d'adjonction (2.8.2) est d'ailleurs familière, chaque fois qu'on a une catégorie $\unsure{H|\mathcal{H}}$ (telle $\mathrm{Cat}$, ou $\mathrm{Hot}_W$ etc) munie d'un foncteur $\pi_0$, \inserted{$\pi_0(X)$} correspondant à l'intuition d'une décomposition d'un objet en ``composantes connexes''. (Je ne vais pas expliciter ici les conditions \inserted{\unsure{générales|génériques|formelles}} dans lesquelles on \unsure{définit|introduit|entend} un tel foncteur.) Par contre, l'adjonction en sens inverse de ces deux mêmes \unsure{foncteurs} est un phénomène tout à fait spécial à des catégories

% archive sheet number: 286
